%% AAMAS 2027 — ACM SigConf format (8 pages + references)
%% Compile: pdflatex paper_aamas && bibtex paper_aamas && pdflatex paper_aamas (×2)

\documentclass[sigconf,review]{acmart}

\usepackage{booktabs}
\usepackage{amsmath}
\usepackage{amssymb}
\usepackage{microtype}
\usepackage{hyperref}

%% ── Metadata ──────────────────────────────────────────────────────────────────
\title{Behavioral Grounding Framework: Empirically Anchored LLM-Based
       Agent Simulations of Synthetic Societies}

\author{Lucas Tourinho}
\email{tourinholucas123@gmail.com}
\affiliation{%
  \institution{TODO: institution}
  \city{TODO: city}
  \country{TODO: country}
}

\begin{document}

\begin{abstract}
Deploying instruction-tuned large language models (LLMs) as agents in
multi-agent economic simulations produces a systematic anomaly: agents
cooperate at rates exceeding 90\%, generating behavioral patterns that
bear no resemblance to observed human populations.  We formalize this
as the \textbf{RLHF Cooperative Bias} and introduce the
\textbf{Behavioral Grounding Framework (BGF)}, a formally specified
agent simulation platform grounded in European Social Survey (ESS Round
11) microdata, to mitigate it.  In a pilot-scale LLM A/B contrast
(single seed, $N{=}50$ agents, $T{=}30$ rounds, Mistral-7B), grounding
reduces cooperation from 96\,\% (Cond.\,A) to 58\,\%, reduces the RLHF
Bias Index ($B_\text{RLHF}$) from 0.71 to 0.25, and reduces the Gini
coefficient from 0.63 to 0.26.  A full-scale $N{=}500$ LLM extension is
pre-registered and launcher-ready.  Cross-cultural validation across six
ESS clusters recovers the cooperation-vs-trust rank ordering perfectly
(Spearman $\rho = +1.000$, $p \approx 0.003$, $n{=}6$).  We release
the complete framework (1,203 automated tests, one-command reproduction)
to enable reproducible research in LLM-driven computational social
science and alignment evaluation.
\end{abstract}

\keywords{LLM agents, agent-based modeling, behavioral grounding,
          RLHF bias, synthetic societies, ESS, multi-agent simulation}

\maketitle

%% ── 1. Introduction ───────────────────────────────────────────────────────────
\section{Introduction}

Agent-based models (ABMs) have long served as tools for studying
emergent social phenomena --- cooperation, inequality, institutional
evolution --- through bottom-up simulation of individual decision-making
\cite{epstein1996,axelrod1984}.  Traditional ABMs rely on hand-crafted
utility functions that are analytically tractable but cannot capture
the linguistic, cultural, and psychological complexity of real human
decision-making.

Recent advances in LLMs offer a compelling alternative: using LLMs as
decision engines for simulated agents (\emph{silicon sampling}
\cite{argyle2023}).  However, LLMs fine-tuned via Reinforcement
Learning from Human Feedback (RLHF) \cite{ouyang2022} produce utopian
societies characterized by frictionless cooperation --- a systematic
distortion we formalize as the \textbf{RLHF Bias Index}
($B_\text{RLHF} = \mathrm{TV}(\pi, \pi_\text{uniform})$).

The \textbf{Behavioral Grounding Framework (BGF)} addresses this
limitation by anchoring LLM agent behavior to empirical microdata from
ESS Round~11.  BGF is formally specified as a tuple
$\mathrm{BGF} = (\mathcal{A}, E, G, P, \Phi, T)$, where
$\Phi: D_\text{ESS} \to \mathit{Profile}$ is an empirical grounding
function that maps survey records to agent profiles, preserving joint
distributions of 15+ sociodemographic attributes.

\paragraph{Research Questions}
\textbf{RQ1}: Does empirical grounding significantly mitigate the RLHF
cooperative bias?
\textbf{RQ2}: How does grounding alter emergent network topology?
\textbf{RQ3}: Can grounded LLM agents reproduce realistic macroeconomic
phenomena?
\textbf{RQ4}: How resilient are grounded societies to adversarial
perturbations and economic shocks?

%% ── 2. Related Work ───────────────────────────────────────────────────────────
\section{Related Work}

% TODO: expand from docs/paper.md §2

\paragraph{Silicon sampling}
Argyle et al.\ \cite{argyle2023} demonstrate that LLMs can simulate
human survey responses with reasonable fidelity, but do not address
multi-agent dynamics or RLHF-induced biases.

\paragraph{Generative Agents}
Park et al.\ \cite{park2023} introduce fictional-persona LLM agents
with episodic memory.  BGF extends this by anchoring personas to
empirical survey distributions and adding dual RAG pipelines for
population-level context injection.

\paragraph{RLHF alignment in social simulation}
% TODO

%% ── 3. Framework ──────────────────────────────────────────────────────────────
\section{Behavioral Grounding Framework}

\subsection{Formal Specification}

$\mathrm{BGF} = (\mathcal{A}, E, G, P, \Phi, T)$ where:
\begin{itemize}
  \item $\mathcal{A}$ — set of agents with profiles $\Phi(d_i)$
        ($d_i \in D_\text{ESS}$)
  \item $E$ — economic environment with action space
        $\{\texttt{work}, \texttt{save}, \texttt{cooperate},
           \texttt{steal}\}$
  \item $G$ — social network (Watts-Strogatz or Erd\H{o}s-R\'enyi)
  \item $P$ — pluggable policy (random / rule-based / LLM / generative)
  \item $\Phi$ — empirical grounding function
        ($D_\text{ESS} \to \mathit{Profile}$)
  \item $T$ — simulation horizon (rounds)
\end{itemize}

\subsection{Behavioral Realism Metric}

% TODO: formal definition from docs/paper.md §3.2

\subsection{RLHF Bias Index}

$B_\text{RLHF} = \mathrm{TV}(\pi_\theta, \pi_\text{uniform})
               = \frac{1}{2}\sum_a |\pi_\theta(a) - \pi_\text{uniform}(a)|$

%% ── 4. Experimental Setup ─────────────────────────────────────────────────────
\section{Experimental Setup}

% TODO: fill from docs/paper.md §4

%% ── 5. Results ────────────────────────────────────────────────────────────────
\section{Results}

\subsection{Pilot A/B Contrast (N=50, T=30, 1 seed)}

% TODO: Table 1 — cooperation rate, B_RLHF, Gini (Cond A vs B)

\subsection{Cross-Cultural Validation}

Spearman $\rho = +1.000$ ($p \approx 0.003$, $n{=}6$, exact permutation
test) across Eastern / Southern / Western / Anglo / Northern / Nordic
ESS clusters.  WVS Wave~7 out-of-sample replication: Pearson $r = +0.977$.

%% ── 6. Discussion ─────────────────────────────────────────────────────────────
\section{Discussion}

% TODO

%% ── 7. Conclusion ─────────────────────────────────────────────────────────────
\section{Conclusion}

We introduced BGF, a formally specified simulation framework that
mitigates the RLHF cooperative bias in multi-agent economic simulations
through empirical grounding in ESS Round~11 microdata.  Pilot results
confirm the central claim: grounding reduces $B_\text{RLHF}$ from 0.71
to 0.25 and cooperation from 96\,\% to 58\,\%.  Cross-cultural
validation independently confirms that BGF recovers real-world trust
rank orderings (Spearman $\rho = 1.0$).  Full-scale GPU runs and Prolific
human evaluation are pre-registered and launcher-ready.

%% ── References ────────────────────────────────────────────────────────────────
\bibliographystyle{ACM-Reference-Format}
\bibliography{paper_aamas}

%% paper_aamas.bib entries (TODO: move to separate .bib file)
%% @article{argyle2023, ...}
%% @article{park2023, ...}
%% @book{epstein1996, ...}
%% @book{axelrod1984, ...}
%% @article{ouyang2022, ...}

\end{document}
